\subsection{Búsqueda binaria}

La búsqueda binaria es un algoritmo para buscar si un elemento específico existe en una secuencia ordenada. El problema que resuelve es entonces:
\begin{itemize}
  \item Entrada: arreglo no vacío ordenado de enteros y número entero $k$
  \item Salida: índice de $k$ en el arreglo o -1 si no está
\end{itemize}

El algoritmo es muy simple: se revisa el elemento del medio del arreglo. Si es $k$, encontramos lo que buscábamos y la salida es el índice del medio; si es menor, solo es sensato seguir buscando en la mitad derecha; si es mayor, solo en la mitad izquierda.

\begin{pseudocode}
def binary_search(nums: arreglo[int], target: int) -> int:
    return binary_search_rec(nums, target, 0, len(nums) - 1)

def binary_search_rec(nums: arreglo[int], target: int, i: int, j: int) -> int:
    |\br{mi}|if j < i:
        return -1|\br{mf}|

    k = |\bxl{over}|i + (j - i) // 2|\bxr{over}|
    if target < nums[k]:
        return binary_search_rec(nums, target, i, k - 1)
    if target > nums[k]:
        return binary_search_rec(nums, target, k + 1, j)

    return k
\end{pseudocode}
\begin{annotations}
  \codebrace[width=6cm]{mi}{mf}{Si \inline{j} está a la izquierda de \inline{i}, los índices se cruzaron, lo que indica que el elemento no está.}
  \codebox[xshift=0.5cm, width=7.5cm]{over}{Una alternativa para calcular la mitad es \inline{(i+j) // 2}, pero arriesga un overflow (en lenguajes de bajo nivel) si la suma es muy grande.}
\end{annotations}

Nótese que en los llamados recursivos se toma desde el índice de la mitad, \inline{k}, \textit{más uno} (o menos uno). Eso cumple dos funciones: primero, que no se vuelva a revisar la mitad; pero más importante aún, permite que a medida que el número no se encuentra y el segmento del arreglo que se contempla es cada vez más pequeño, de forma que el índice de la mitad se acerca más y más al extremo, eventualmente se crucen y la recursión llegue al caso base. Si solo se tomara el índice de la mitad, jamás se cruzaría con el extremo y la recursión no terminaría. Y no se puede detener cuando son iguales, porque ese caso de un único elemento sí se quiere revisar.

Como cada vez que se parte el problema se reduce a la mitad, para determinar la complejidad se quiere saber cuántas veces se debe dividir la entrada de tamaño $n$ entre $2$ hasta que se llega a un problema de tamaño 1. Para conocer esa $x$ cantidad de veces, se formula:

\begin{gather*}
  n \underbrace{{} \cdot \frac{1}{2} \cdot \frac{1}{2} \cdots \frac{1}{2}}_{x \ \text{veces}} = 1 \\
  \frac{n}{2^x} = 1
\end{gather*}

Despejando,

\begin{gather*}
  2^x = n \\
  x = \log_2(n)
\end{gather*}
Es decir, el número de divisiones, que es el número de pasos, crece logarítmicamente con respecto al tamaño de la entrada.

En este algoritmo, solo se realizan las etapas de dividir y conquistar; no es necesario combinar las soluciones. Eso facilita realizar una implementación iterativa del algoritmo:
\begin{pseudocode}
def binary_search_iter(nums: arreglo[int], target: int) -> int:
    i = 0
    j = len(nums) - 1
    
    idx = -1

    while i <= j and idx == -1:
        k = i + (j-i) // 2
        if target < nums[k]:
            j = k - 1
        elif target > nums[k]:
            i = k + 1
        else:
            idx = k
        
    return idx
\end{pseudocode}


\practicar[leetcode]{Binary Search}{https://leetcode.com/problems/binary-search/}


